%! Conic Sections — Unit 8, Lesson 8.5 — Guided Notes
\documentclass[10pt]{article}
\usepackage{precalculus-article}
\usepackage{precalculus-boxes}
\newcommand{\TallMath}[1]{$\displaystyle #1\rule[-1.4em]{0pt}{3.2em}$}

\begin{document}

\pageheader{Unit 8, Lesson 8.5}{Guided Notes}
\namedateperiod

\vspace{0.10in}

%----------------------------------------------------------------------
% Objectives
%----------------------------------------------------------------------
\begin{objectivebox}
\textbf{I will be able to\ldots}
\begin{enumerate}[leftmargin=1.5em, itemsep=4pt]
  \item \textbf{LO 4.6.A} --- Write the standard form equation of a parabola,
    ellipse (including circles), or hyperbola given key features; and identify
    those features from a given equation. \writeline
\end{enumerate}
\end{objectivebox}

\vspace{0.08in}

%----------------------------------------------------------------------
% Vocabulary
%----------------------------------------------------------------------
\begin{vocabbox}
\begin{description}[leftmargin=0pt, itemsep=6pt]
  \item[\termblanklong{conic section}]
    A curve formed by the intersection of a plane and a double cone:
    a \underline{\hspace{2cm}}, an \underline{\hspace{2cm}}, or a
    \underline{\hspace{2cm}}.

  \item[\termblanklong{vertex (parabola)}]
    The point $(h, k)$ in the standard form $y - k = a(x-h)^2$ or
    $x - h = a(y-k)^2$. It is the \underline{\hspace{2.5cm}} point of the parabola.

  \item[\termblanklong{center (ellipse / hyperbola)}]
    The point $(h, k)$ at the \underline{\hspace{2.5cm}} of an ellipse or hyperbola;
    read from $h$ and $k$ in the standard form.

  \item[\termblanklong{semi-major / semi-minor axis}]
    For an ellipse: the larger radius is $a$ (\underline{\hspace{2cm}}) and the
    smaller is $b$ (\underline{\hspace{2cm}}).

  \item[\termblanklong{asymptote (hyperbola)}]
    Lines the hyperbola approaches but never reaches:
    $y - k = \pm\dfrac{\;\;}{\;\;}(x - h)$. \quad Fill in $b$ over $a$.
\end{description}
\end{vocabbox}

\vspace{0.08in}

%----------------------------------------------------------------------
% Hook
%----------------------------------------------------------------------
\begin{hookbox}
\textbf{Entry Question:} Planetary orbits, telescope mirrors, and navigation
systems all use conic sections.

\vspace{4pt}
\textbf{Q1.} Name one real-world application of a parabola and one of an ellipse.
\writelines{2}

\textbf{Q2.} What do you predict all three conic shapes have in common
algebraically?
\writelines{2}
\end{hookbox}

\vspace{0.08in}

%----------------------------------------------------------------------
% Section 1 — Parabolas
%----------------------------------------------------------------------
\begin{notesbox}{Section 1 --- Parabolas (LO 4.6.A, EK 4.6.A.1)}

\textbf{Standard forms:}

\vspace{4pt}
\renewcommand{\arraystretch}{1.8}
\begin{tabularx}{\linewidth}{@{} p{0.38\linewidth} p{0.30\linewidth} X @{}}
\textbf{Form} & \textbf{Axis} & \textbf{Direction} \\
\midrule
$y - k = a(x - h)^2$ & vertical & $a > 0$: opens \underline{\hspace{1.2cm}};
  $a < 0$: opens \underline{\hspace{1.2cm}} \\[4pt]
$x - h = a(y - k)^2$ & horizontal & $a > 0$: opens \underline{\hspace{1.2cm}};
  $a < 0$: opens \underline{\hspace{1.2cm}} \\
\end{tabularx}

\vspace{8pt}
\textbf{Worked Example:}\quad $y + 3 = 2(x - 1)^2$

Rewrite: $y - (\underline{\phantom{-3}}) = 2(x - \underline{\phantom{1}})^2$

Vertex: \blank{3cm} \quad $a = \blank{1.2cm}$, so opens: \blank{2cm}

\vspace{10pt}
\textbf{Student Practice:}\quad $y - 5 = -(x + 2)^2$

Vertex: \blank{3cm} \quad $a = \blank{1.2cm}$, so opens: \blank{2cm}

\end{notesbox}

\vspace{0.08in}

%----------------------------------------------------------------------
% Section 2 — Ellipses and Circles
%----------------------------------------------------------------------
\begin{notesbox}{Section 2 --- Ellipses and Circles (LO 4.6.A, EK 4.6.A.2)}

\textbf{Standard form of an ellipse:}
\[
  \frac{(x-h)^2}{a^2} + \frac{(y-k)^2}{b^2} = 1
\]

\textbf{Worked Example:}\quad
$\dfrac{(x+1)^2}{16} + \dfrac{(y-3)^2}{9} = 1$

\vspace{6pt}
Center: \blank{3cm} \quad
$a^2 = \blank{1.2cm} \Rightarrow a = \blank{1.2cm}$ (horizontal) \quad
$b^2 = \blank{1.2cm} \Rightarrow b = \blank{1.2cm}$ (vertical)

\vspace{10pt}
\textbf{Circle (special case):}\quad When $a = b = r$, the equation becomes

\[
  (x - h)^2 + (y - k)^2 = r^2
\]

\textbf{Student Practice:}\quad
$\dfrac{(x-3)^2}{25} + \dfrac{(y+1)^2}{4} = 1$

Center: \blank{3cm} \quad $a = \blank{1.2cm}$ \quad $b = \blank{1.2cm}$

\end{notesbox}

\vspace{0.08in}

%----------------------------------------------------------------------
% Section 3 — Hyperbolas
%----------------------------------------------------------------------
\begin{notesbox}{Section 3 --- Hyperbolas (LO 4.6.A, EK 4.6.A.3)}

\textbf{Standard forms:}

\vspace{4pt}
\renewcommand{\arraystretch}{1.8}
\begin{tabularx}{\linewidth}{@{} p{0.44\linewidth} X @{}}
\textbf{Form} & \textbf{Opens} \\
\midrule
$\dfrac{(x-h)^2}{a^2} - \dfrac{(y-k)^2}{b^2} = 1$ &
  left and right ($x$-term positive) \\[4pt]
$\dfrac{(y-k)^2}{b^2} - \dfrac{(x-h)^2}{a^2} = 1$ &
  up and down ($y$-term positive) \\
\end{tabularx}

\vspace{8pt}
\textbf{Asymptotes:}\quad $y - k = \pm\dfrac{b}{a}(x - h)$

\vspace{10pt}
\textbf{Worked Example:}\quad
$\dfrac{(x-2)^2}{9} - \dfrac{(y+1)^2}{4} = 1$

Center: \blank{3cm} \quad $a = \blank{1.2cm}$ \quad $b = \blank{1.2cm}$ \quad
Opens: \blank{2cm}

Asymptotes: $y + 1 = \pm\blank{2.5cm}(x - 2)$

\vspace{10pt}
\textbf{Student Practice:}\quad
$\dfrac{(y+3)^2}{16} - \dfrac{(x-1)^2}{9} = 1$

Center: \blank{3cm} \quad $a = \blank{1.2cm}$ \quad $b = \blank{1.2cm}$ \quad
Opens: \blank{2cm}

Asymptotes: $y + 3 = \pm\blank{2.5cm}(x - 1)$

\end{notesbox}

\end{document}
